\section{Broader Impacts and Asset Notes}
\label{app:responsible-assets}

Our method is a schedule-calibration method, separate from model or dataset
release. Its
positive impact, when reuse is possible, is reduced generation error at a fixed
NFE or fewer failed generations for a fixed model. Its main negative impact is
that it can also make existing generative-image systems cheaper or easier to
tune. Such systems inherit risks from their base models, including harmful
content, social bias, mis- and disinformation, impersonation, and copyright or
license misuse. The matched SD1.5 and SDXL results in this paper are reported
as automated-metric boundary cases, leaving broad text-to-image quality and
safer generation unsupported.

It also adds calibration compute. In settings where a calibrated
schedule is used only once, total compute can increase sharply. Any
deployment claim should therefore account for schedule reuse, safety filtering,
and model-license restrictions. This work uses no human subjects, private data,
or deployed system. Complete local paths for the assets below are recorded in
\texttt{configs/assets\_manifest.yaml}.

\begin{table}[t]
\centering
\caption{Current asset audit. ``Pending'' marks items that need final release
documentation before submission.}
\label{tab:asset-audit}
\scriptsize
\setlength{\tabcolsep}{3pt}
\begin{tabular}{p{0.18\linewidth}p{0.16\linewidth}p{0.30\linewidth}p{0.22\linewidth}}
\toprule
Asset & Role & Source/license status & Local record \\
\midrule
CIFAR-10 & FID benchmark and reference statistics &
Created by Krizhevsky, Nair, and Hinton; UCI lists CC BY 4.0; cite the
technical report \citep{krizhevsky2009learning}. &
CIFAR config and FID-stat asset \\
PNDM CIFAR checkpoint & Controlled sampler backend &
The official PNDM repository is Apache-2.0 and links the checkpoint/statistics
folder; cite \citet{liu2022pndm}. The release manifest records the local CIFAR
checkpoint and FID-stat hashes, but the archive should not bundle checkpoint
files without final redistribution review. &
PNDM CIFAR checkpoint asset \\
Stable Diffusion 1.5 & Text-to-image backend &
Cached model card lists \texttt{creativeml-openrail-m}. &
SD1.5 cached model card \\
SDXL & Text-to-image backend &
Cached model card and upstream terms need final author review. &
SDXL cached model record \\
AYS 10-step schedule & Published offline schedule baseline &
The official AYS quickstart lists the 10-step values; the release manifest
records local hashes. The anonymous package omits raw numeric schedule files
and retains AYS through citations, hashes, rendered figures, and metric rows. &
AYS SD1.5 and SDXL 10-step bundles \\
Diffusers & Text-to-image execution library &
Repository license is Apache-2.0. &
Diffusers vendored checkout \\
ImageReward & Text-to-image metric &
Model card and repository license are Apache-2.0. &
ImageReward checkout \\
Prompt set & Matched text-to-image prompts &
Project-local 50-prompt asset; the paper-facing manifest records hashes, while
the release license still needs final audit. &
SD prompt asset \\
\bottomrule
\end{tabular}
\end{table}

The new assets produced by this paper are experiment configs, schedule
metadata, cleaned CSVs, tables, plots, and manuscript text under
\texttt{paper/}. Draft release metadata now lives under
\texttt{paper/release/}, including an anonymized artifact manifest, license
status note, and generated-sample policy. Before release, the authors still
need to choose final licenses, finish the prompt redistribution audit, and
prepare an external anonymous archive or repository. The current release policy
omits raw AYS numeric schedule files from the anonymous package.
